\documentclass[10pt,twocolumn,letterpaper]{article}

\usepackage[letterpaper,margin=0.85in]{geometry}
\usepackage{titling}
\setlength{\columnsep}{0.3in}
\setlength{\droptitle}{-3em}
\renewcommand{\abstractname}{\large Abstract}
\usepackage{cite}
\usepackage{amsmath,amssymb,amsfonts}
\usepackage{graphicx}
\usepackage{textcomp}
\usepackage{xcolor}
\usepackage{booktabs}
\usepackage{multirow}
\usepackage{url}
\usepackage[breaklinks=true,hidelinks]{hyperref}
\usepackage{caption}
\captionsetup{font=footnotesize}

\title{Feature-Augmented Graph Neural Networks for\\
Out-of-Distribution Detection of Collusive\\
Market Manipulation: A Multi-Cohort Study}

\author{Sumit Sontakke\\
\small Department of Computer Science, PES University, Bengaluru, India\\
\small \texttt{pes2pge24ds111@stu.pes.edu}}

\begin{document}
\maketitle

\begin{abstract}
Surveillance systems that flag collusive trading must do more than fit
the data they were trained on; they must hold up when the parameters,
seeds, and even the data-generating process shift. We report a
multi-stage study of a two-layer GraphSAGE manipulation detector
augmented with six taxonomy-grounded per-trader features, evaluated
across three increasingly difficult generalisation regimes.
On a fifty-run out-of-distribution cohort with parameter-disjoint
manipulation scenarios, the feature-augmented model lifts edge-level
ROC AUC from \(0.638\) (two topology features) to \(0.984\) and trader
recall from \(43\%\) to \(83\%\) at \(74\%\) purity. The same features,
fed instead into a gradient-boosted second-stage classifier on top of
the baseline graph network, reach pooled trader AUC \(0.968\) under
leave-one-run-out cross-validation and clear the \(99\%\)-precision bar
at \(62\%\) recall.
On a fifteen-run cross-generator cohort whose manipulators are produced
by a different agent-based simulator entirely, the augmented network
retains edge AUC \(0.842\) while the two-feature baseline collapses to
chance \((0.518)\), demonstrating that the engineered features carry
manipulation-essential signal across simulation regimes.
A leave-one-family-out test, however, reveals an unexpected contrast:
the second-stage classifier maintains AUC between \(0.91\) and \(0.99\)
on the held-out family, while the end-to-end network collapses to chance
in every fold. The simpler bolt-on stack therefore generalises across
manipulation styles better than the more expressive end-to-end model.
We discuss the implications for surveillance system design and outline
the open questions for cross-distribution operating-point calibration.
\end{abstract}

\vspace{0.5em}\noindent\textbf{Keywords:} 
Market surveillance, collusive manipulation, graph neural networks,
out-of-distribution generalisation, agent-based simulation,
feature ablation, leave-one-family-out validation.


\section{Introduction}
\label{sec:intro}

Detecting collusive manipulation in equity markets is a problem whose
character has changed substantially over the past decade. Standalone
abusive behaviour by a single trader can often be caught with simple
rules on order volume and price impact, but coordinated activity by a
ring of accounts, a clique acting through proxy fronts, or a group
deliberately splitting trades across counterparties leaves only a
relational signature. Each individual trader may look unremarkable;
the conspiracy is visible only when one looks at \emph{who trades with
whom, when, and in what coordinated pattern}. This makes the detection
task naturally graph-shaped and motivates the use of graph neural
networks (GNNs), which can encode both per-trader features and the
structure of inter-trader interactions in a unified representation.

Yet expressive models are not, by themselves, enough. A surveillance
system is judged not on its training distribution but on its operating
distribution, and the operating distribution is almost always different.
Real markets shift by season, by ticker, by regulator-induced
microstructure changes; even when the underlying generator is held
fixed, the parameter envelope of any specific period is unlikely to
match the model's training period. Worse, the synthetic cohorts on which
academic detectors are typically validated come from a single simulator,
and a model that succeeds within that simulator's parameter range may
still fail when exposed to a structurally different generator. Robust
deployment therefore requires multiple, increasingly demanding
generalisation tests rather than a single in-distribution accuracy
number.

This paper presents a study that does exactly that. We build a two-layer
GraphSAGE model on a calibrated synthetic market substrate and study its
behaviour under three nested out-of-distribution regimes: a
parameter-disjoint test cohort drawn from the same generator, a
cross-generator cohort whose manipulators come from a different
agent-based simulator (ABIDES \cite{byrd2019abides}), and a
leave-one-family-out evaluation in which one of the three manipulation
families (clique, ring, mixed) is held out from training entirely. To
contextualise the model-only result, we also evaluate a second-stage
gradient-boosted classifier that consumes the GraphSAGE output together
with six per-trader engineered features designed to capture
manipulation-essential structure.

Our contributions are fourfold. First, we present a six-feature
engineered representation of trader-level manipulation signatures and
show that several of its components carry standalone discriminative
power comparable to or exceeding that of the baseline graph network.
Second, we show that simply concatenating these six features onto the
GNN node vector and retraining from scratch lifts the model's
within-generator OOD edge AUC from \(0.638\) to \(0.984\). Third, we
report a cross-generator evaluation on ABIDES manipulator agents in
which the augmented model holds AUC \(0.842\) (versus chance for the
unaugmented baseline), confirming that the engineered features carry
signal across simulators. Fourth, and most surprisingly, we report a
leave-one-family-out evaluation in which the simple second-stage
classifier remains accurate (AUC \(0.91\)--\(0.99\)) while the
end-to-end augmented network collapses to chance. We discuss why this
appears to be the case and what it implies for surveillance system
design.

The remainder of this paper is organised as follows. Section
\ref{sec:related} surveys related work in manipulation detection and
graph learning. Section \ref{sec:problem} formalises the problem.
Section \ref{sec:data} describes the synthetic cohorts used. Section
\ref{sec:methods} describes the four detector architectures we
compare. Section \ref{sec:results} reports the within-generator,
cross-generator, and family-disjoint results. Section \ref{sec:disc}
discusses what the comparison implies. Sections \ref{sec:lim} and
\ref{sec:concl} cover limitations and conclude.

\section{Related Work}
\label{sec:related}

\textbf{Market manipulation detection} has a long literature, much of
it rule- or threshold-driven. Early work focused on layering, spoofing,
and pump-and-dump patterns visible in order-book microstructure.
Supervised classifiers later appeared,
treating manipulation detection as a per-order or per-trader
classification problem with hand-engineered statistics; Cao and
co-authors, for instance, applied particle-swarm-tuned learners to
labelled real and synthetic exchange data \cite{cao2014pso}. The shift
from per-order to relational reasoning has been driven by the
observation that genuinely collusive activity tends to evade
trader-marginal classifiers: any single account looks legitimate.
\textbf{Graph-based fraud and anomaly detection} has emerged as a
natural response. Recent work has applied graph neural networks to
financial-fraud datasets and demonstrated meaningful lifts over
non-relational baselines, though most
benchmarks remain within-distribution.

\textbf{Graph neural networks for inductive learning.} We adopt
GraphSAGE \cite{hamilton2017graphsage} as the backbone because it
supports inductive node-representation learning, which matters in the
surveillance setting where deployment graphs differ from training
graphs in both the identity of nodes and the size of the graph. Other
attention-based architectures such as GAT \cite{velickovic2018gat} and
spectral methods such as GCN \cite{kipf2017gcn} are widely studied
alternatives, but their inductive properties on dynamic, partially
observed market graphs are less well characterised.

\textbf{Imbalanced classification} pervades surveillance: only a small
fraction of any cohort is manipulative. We use focal loss
\cite{lin2017focal} for the GNN's edge-prediction head, which down-weights
easy negatives and concentrates the gradient on hard positives. For the
second-stage classifier we use gradient-boosted trees
\cite{friedman2001gradient}, which handle modest training set sizes and
heterogeneous features well without extensive tuning.

\textbf{Out-of-distribution evaluation.} The need to evaluate
classifiers on data drawn from a distinct distribution has been
emphasised in image classification \cite{recht2019imagenet} and made
systematic in domain-generalisation benchmarks such as WILDS
\cite{koh2021wilds}. The natural analogue in the surveillance setting
is to vary not only the seeds and parameters but also the underlying
data-generating process. To that end, we use ABIDES
\cite{byrd2019abides}, a discrete-event agent-based market simulator,
to produce a held-out cohort whose generator differs from our primary
synthesiser.

\textbf{Calibrated synthetic markets.} The baseline synthesiser used in
this work follows the stylised facts of asset returns set out by Cont
\cite{cont2001stylized}, including heavy-tailed return distributions,
volatility clustering, and the absence of linear autocorrelation in
returns. Calibration to a real exchange (the Indian NSE bhavcopy time
series) was performed in an earlier project phase to ensure the
synthetic market is at least first-moment-faithful to real conditions.

\section{Problem Formulation}
\label{sec:problem}

We consider a daily slice of an equity market populated by
\(N\)~traders, each of whom emits orders that may or may not match
against the limit order book. Trades resulting from matched orders
form a directed weighted graph \(G=(V,E)\) in which each trader is a
node and an edge from \(a\) to \(b\) records that \(a\) sold to \(b\)
some quantity over the slice. A subset \(M\subset V\) of traders is
deemed manipulative; in the synthetic regime considered here this
labelling is exact, with each \(M\) belonging to one of three
manipulation \emph{families}:

\begin{itemize}
\item \textbf{Clique}: a small set of accounts that coordinate to push
or absorb price, typically routed through one or two shared front
counterparties.
\item \textbf{Ring}: a larger set of accounts that pass positions
around a closed cycle, creating heavy fictitious volume.
\item \textbf{Mixed}: a scenario combining elements of both, in which
some accounts participate in clique behaviour and others in ring
behaviour during the same slice.
\end{itemize}

The detector's task is to output a score
\(\hat{y}\colon V\to[0,1]\) ranking each trader by manipulator
probability, ideally such that the score's ordering is stable under
seed, parameter and generator shifts. We evaluate via three metrics:
edge-level ROC AUC (informative because the underlying classifier
emits edge probabilities first), per-family trader recall at a locked
threshold, and per-family trader purity (precision among flagged
traders) at the same threshold.

\section{Data and Cohort Design}
\label{sec:data}

\subsection{Training Cohort}
The training cohort consists of five days of synthetic market
sessions; each day comprises fifty independent runs sampled from a
parameter envelope deliberately wider than the typical evaluation
suite. Each run contains approximately 500 traders, of which
\(0.05\%\)\,-\,\(2\%\) are manipulators belonging to a randomly drawn
family. Manipulation duration is sampled between 10 and 60 minutes
within the trading session. The five-day cohort design supports a
continual warm-start training protocol in which the model retains
weights from one day's training to seed the next, with smaller per-day
epoch budgets after the initial day to limit catastrophic forgetting.

\subsection{Within-Generator OOD Cohort}
The primary out-of-distribution evaluation cohort consists of fifty
runs drawn from a parameter envelope strictly outside the training
distribution: session duration \(60\)--\(90\) minutes, manipulation
ratio \(2\)--\(5\%\), core manipulator count between \(16\) and
\(25\). Random seeds are held out from the training distribution.
Approximately one-third of the runs are clique scenarios, one-third
ring, and one-third mixed.

\subsection{Cross-Generator Cohort}
To stress the model further, we generated an additional eighteen runs
using ABIDES, an agent-based discrete-event market simulator
distinct from the primary synthesiser. Three pilot runs (one per
family) verified that the data schema matched, and fifteen further
runs (five seeds across three families, all at a held-out calibration
date) form the statistical-power expansion. Each ABIDES run contains
approximately 125 active traders and six manipulator agents per run.
The ABIDES manipulator agents---\texttt{CollusiveCliqueAgent},
\texttt{RingTraderAgent}, and \texttt{FrontAccountAgent}---are
implementation-level approximations of the same three taxonomies used
in the primary synthesiser, but the underlying microstructure (limit
order book dynamics, message ordering, exchange agent behaviour)
differs substantially.


\subsection{Dataset Statistics}
\label{sec:stats}
Table~\ref{tab:cohorts} summarises the three cohorts used in this
study. The training cohort is the largest by run count; the
within-generator OOD cohort is comparable in per-run size; the
cross-generator ABIDES cohort is smaller per run because the ABIDES
agent population is configured around \(125\) active traders rather
than the \(500\) used in the primary synthesiser, but it remains
large enough to expose distribution shift without becoming trivial.

\begin{table}[t]
\caption{Cohort summary.}
\label{tab:cohorts}
\centering
\footnotesize
\begin{tabular}{lccc}
\toprule
Cohort & Runs & Traders/run & Manip./run \\
\midrule
Training (5-day) & 250 & $\sim$500 & $\sim$10--20 \\
Within-generator OOD & 50 & $\sim$500 & $\sim$16--25 \\
ABIDES pilot & 3 & $\sim$125 & 6 \\
ABIDES expanded & 15 & $\sim$125 & 6 \\
\bottomrule
\end{tabular}
\end{table}

\section{Detector Architectures}
\label{sec:methods}

\subsection{Baseline GraphSAGE (v1)}
\label{sec:v1}
Our baseline detector is a two-layer GraphSAGE network operating on the
per-run directed trade graph. Each layer applies a
\texttt{SAGEConv} aggregator (we use the \texttt{max}-pool variant) with
hidden dimensions \((256,128,64)\) and a dropout rate of \(0.3\). An
MLP head consumes the concatenated source-destination embeddings and
edge-level features---lead-lag correlation, signed imbalance, traded
volume, interaction count, pre- and post-window drift, and VWAP
imbalance---and emits a per-edge manipulation probability. The model is
trained with focal loss
\((\alpha{=}0.85,\gamma{=}2.0)\) and a custom early-stopping signal
based on the per-family trader recall on a held-out validation set
rather than the more conventional validation cross-entropy. This
training protocol was chosen for its alignment with the surveillance
operating metric.

Each graph's node feature vector consists of just two features:
the trader's total executed quantity and the trader's number of
distinct counterparties. The edge feature vector contains the eight
quantities listed above, computed over a one-minute bucket.

After training, edge probabilities are projected to trader scores via
the rule
\(\hat{y}(v) = 0.7\max_e p(e) + 0.3\,\mathrm{mean}_{e\in\mathrm{top}_3}\, p(e)\),
which gave the best per-family balance in earlier ablations.

\subsection{Engineered Manipulation-Signature Features}
\label{sec:features}
We define six per-trader features each scoped to that trader's
busiest five-minute order window. Below, \(O_v\) denotes the set of
orders placed by trader \(v\) during the session, \(B_v\subseteq O_v\)
denotes their burst window, and \(T_v\) denotes the trades to which
\(v\) was a party during \(B_v\).

\begin{enumerate}
\item \textbf{Burst concentration}
\(\phi_1(v) = |B_v|/|O_v|\).
\item \textbf{Side entropy in burst}
\(\phi_2(v) = -p\log_2 p - (1-p)\log_2(1-p)\)
where \(p\) is the buy fraction within \(B_v\).
\item \textbf{Counterparty HHI in burst}
\(\phi_3(v) = \sum_k s_k^2\),
where \(s_k\) is the share of \(T_v\) attributable to counterparty \(k\).
\item \textbf{Order quantity coefficient of variation}
\(\phi_4(v) = \sigma(q)/\mu(q)\)
where \(q\) is the vector of order quantities of \(v\).
\item \textbf{Top-partner trade share}
\(\phi_5(v) = \max_k s_k\), using the same shares as \(\phi_3\).
\item \textbf{Co-active top-quantile count}
\(\phi_6(v)\): the number of other top-activity-quantile traders whose
burst window overlaps that of \(v\).
\end{enumerate}

The standalone discriminative power of each feature on the primary OOD
cohort is given in Table~\ref{tab:standalone}. The strongest single
feature is \(\phi_3\) (counterparty HHI) at AUC \(0.800\); the weakest
is \(\phi_2\) (side entropy) at \(0.550\). Notably, three of the six
exceed the baseline GraphSAGE's pooled trader AUC of \(0.793\),
indicating that the features alone carry signal comparable to the
end-to-end network.

\begin{table}[t]
\caption{Standalone ROC AUC of engineered features on the primary OOD cohort.}
\label{tab:standalone}
\centering
\begin{tabular}{lc}
\toprule
Feature & AUC \\
\midrule
Counterparty HHI in burst ($\phi_3$) & 0.800 \\
Top-partner trade share ($\phi_5$) & 0.775 \\
Co-active top-quantile count ($\phi_6$) & 0.744 \\
Burst concentration ($\phi_1$) & 0.704 \\
Order qty coefficient of variation ($\phi_4$) & 0.666 \\
Side entropy in burst ($\phi_2$) & 0.550 \\
\midrule
Baseline GraphSAGE v1 trader score (reference) & 0.793 \\
\bottomrule
\end{tabular}
\end{table}

\subsection{Tier-2 Stacked Classifier}
\label{sec:tier2}
Our first augmented detector is a stacked classifier. We take the
baseline GraphSAGE's per-trader score as input \(x_0\), concatenate the
six engineered features \((\phi_1,\dots,\phi_6)\) to form a seven-element
input vector, and train a gradient-boosted classifier
\((n=120,\text{depth}{=}3,\eta{=}0.08,\text{subsample}{=}0.85)\) on
this representation under leave-one-run-out cross-validation. The
classifier emits a trader-level manipulator probability that combines
the graph network's view (\(x_0\)) with the trader-marginal features.

\subsection{Feature-Augmented GraphSAGE (v4)}
\label{sec:v4}
The second augmented detector folds the engineered features into the
graph network's node representation instead of using them as a
downstream classifier. The node feature vector grows from two
dimensions to eight: \([\text{volume},\text{counterparties},
\phi_1,\dots,\phi_6]\). The remaining architecture, training schedule,
and loss are held fixed, so the only variable between the baseline
network and v4 is the node feature set.

This means v4 cannot warm-start from the baseline checkpoint---the
first \texttt{SAGEConv} layer's weight matrix changes shape. v4
therefore trains from scratch on the same five-day cohort.


\subsection{Trader-Level Projection Rule}
\label{sec:proj}
All four detectors emit per-trader scores. For the two graph networks,
which natively produce per-edge probabilities, we project to traders
via a weighted blend of each trader's maximum incident-edge probability
and the mean of their top three. Formally, for trader \(v\) with
incident-edge probabilities \(\{p_e\}_{e\in E(v)}\):
\begin{equation}
\hat{y}(v) \;=\; 0.7\,\max_{e\in E(v)} p_e \;+\; 0.3\,
\frac{1}{3}\!\!\sum_{e\in\mathrm{top}_3(v)}\!\! p_e .
\label{eq:proj}
\end{equation}
The \(0.7\)/\(0.3\) split was selected in an earlier project phase by
maximising per-family trader recall on a held-out validation cohort.
A pure-\texttt{max} projection over-weights single noisy edges and
hurts purity; a pure-mean projection dilutes the strong signal from
clear manipulators. Equation~\eqref{eq:proj} balances these tendencies.

\subsection{Training Protocol}
\label{sec:training}
The training protocol is summarised in Algorithm~\ref{alg:train}.
Five days of cohort data are processed sequentially. Day~0 is trained
from scratch with the full epoch budget; days~1--4 warm-start from the
previous day's checkpoint with a smaller epoch budget to limit
catastrophic forgetting. Early stopping is governed by per-family
trader recall on an in-day validation split rather than by validation
loss, in keeping with the surveillance operating metric.

\begin{figure}[t]
\small
\noindent\textbf{Algorithm}: Continual warm-start GraphSAGE training.
\begin{enumerate}
\item \textbf{Input:} cohort days $D_0,\dots,D_4$; node-feature spec $\Phi$
\item \textbf{Output:} per-day checkpoints $\theta_0,\dots,\theta_4$
\item Initialise model $\theta\leftarrow$ random
\item \textbf{for} $d=0,\dots,4$ \textbf{do}
  \item $G_d \gets \{\text{build\_graph}(r,\Phi)\;|\;r \in D_d\}$
  \item $(G_d^{\text{tr}}, G_d^{\text{val}}) \gets$ split$(G_d, p=0.10)$
  \item epochs $\gets 80$ if $d=0$ else $25$
  \item Train $\theta$ with focal loss on $G_d^{\text{tr}}$,
    early-stop on val-trader-recall($G_d^{\text{val}}$),
    patience $6$
  \item $\theta_d \gets \theta$; persist checkpoint
\item \textbf{end for}
\end{enumerate}
\label{alg:train}
\end{figure}


\subsection{Feature Computation}
\label{sec:featcomp}
Each of the six engineered features in \S\ref{sec:features} is computed
from order- and trade-level events scoped to a per-trader \(5\)-minute
burst window. Algorithm~\ref{alg:feat} sketches the procedure. The
burst window is selected as the contiguous \(5\)-minute slice with the
highest order count for that trader. Features are computed per run and
cached, so repeated experiments incur only one feature-computation pass
per cohort regardless of model.

\begin{figure}[t]
\small
\noindent\textbf{Algorithm}: Per-trader engineered feature extraction.
\begin{enumerate}
\item \textbf{Input:} orders $O$, trades $T$ for one run; trader $v$
\item \textbf{Output:} $\phi_1,\dots,\phi_6$ for $v$
\item $O_v \gets$ orders of $v$ in $O$
\item $B_v \gets$ argmax over $5$-min windows of $|O_v \cap \text{window}|$
\item $T_v \gets$ trades involving $v$ during $B_v$
\item $\phi_1 \gets |B_v| / |O_v|$
\item $p \gets$ buy-fraction in $B_v$;
  $\phi_2 \gets -p\log_2 p - (1-p)\log_2(1-p)$
\item $s_k \gets$ counterparty-$k$ share of $T_v$;
  $\phi_3 \gets \sum_k s_k^2$;
  $\phi_5 \gets \max_k s_k$
\item $\phi_4 \gets \sigma(q)/\mu(q)$ for order quantities $q$ in $O_v$
\item $A \gets$ top-quartile-activity traders active during $B_v$;
  $\phi_6 \gets |A \setminus \{v\}|$
\end{enumerate}
\label{alg:feat}
\end{figure}


\section{Implementation and Reproducibility}
\label{sec:impl}
\subsection{Hardware and Software Stack}
All graph network training and inference runs on an NVIDIA RTX
\(5080\) workstation card (sm\_120 architecture, \(16\)\,GB VRAM). The
training pipeline is containerised with Docker; the GraphSAGE trainer
uses PyTorch \(2.x\) and PyTorch Geometric. The stacked classifier
uses scikit-learn's \texttt{GradientBoostingClassifier}. Per-run
graph construction is pure NumPy/pandas and is independent of the
deep-learning stack, which keeps the feature pipeline CPU-portable.

\subsection{Hyperparameters}
Table~\ref{tab:hparams} lists the hyperparameters held constant across
all comparisons. The graph network's loss configuration was selected
in an earlier ablation in which a softer focal
(\(\gamma\!=\!1.0,\alpha\!=\!0.75\)) and longer patience were both
empirically worse than the baseline. The boosted classifier's settings
are scikit-learn defaults adjusted for the dataset size; we did not
perform a hyperparameter search beyond a small initial grid for the
number of estimators.

\begin{table}[t]
\caption{Hyperparameters held constant across all comparisons.}
\label{tab:hparams}
\centering
\footnotesize
\begin{tabular}{lll}
\toprule
Component & Parameter & Value \\
\midrule
\multirow{6}{*}{GraphSAGE} & layers & 2 \\
 & hidden dims & 256, 128, 64 \\
 & aggregator & max \\
 & dropout & 0.3 \\
 & loss & focal $(\alpha{=}0.85,\gamma{=}2.0)$ \\
 & early-stop patience & 6 (val trader recall) \\
\midrule
\multirow{4}{*}{Tier-2 GBM} & n\_estimators & 120 \\
 & max\_depth & 3 \\
 & learning\_rate & 0.08 \\
 & subsample & 0.85 \\
\midrule
Both & random seed & 42 \\
\bottomrule
\end{tabular}
\end{table}

\subsection{Runtime}
A single five-day continual training run for the augmented network
takes approximately \(170\)~minutes on the workstation card. Feature
caching for the OOD cohort completes in under two minutes on CPU. The
boosted classifier under leave-one-run-out cross-validation
(\(50\) folds, \(\approx\!10{,}000\) trader-rows) completes in under
\(45\)~seconds on a single CPU core. Inference on the \(50\)-run OOD
cohort for either model takes well under a second on GPU. These
modest runtimes mean the entire ablation matrix in this paper was
produced from disk-cached graphs and features in a few hours of
wall-clock time, supporting frequent iteration.

\section{Experiments}
\label{sec:results}

\subsection{Hyperparameter Ablation}
We first verified that the baseline GraphSAGE's hyperparameters were
near-optimal before declaring the model the ceiling. Three
single-knob variants were retrained: v1 (the baseline; focal
\(\gamma=2.0,\alpha=0.85\), patience \(=6\)), v2
(\(\gamma=1.0,\alpha=0.75\), patience \(=12\)), and v3
(\(\gamma=2.0,\alpha=0.85\), patience \(=12\)). Trader-level AUCs were
\(0.612\), \(0.486\), and \(0.083\) respectively. The patience knob
alone collapsed the model. We concluded that hyperparameter tuning had
hit a ceiling and that further improvement would require either richer
features or richer model structure.

\subsection{Within-Generator OOD Result}
\label{sec:result-ood}
The within-generator out-of-distribution result is the primary headline
of the study. Table~\ref{tab:ood} compares the baseline network against
its feature-augmented retraining on the fifty-run OOD cohort. Edge AUC
climbs from \(0.638\) to \(0.984\); per-family trader recall climbs
from \(0.41\)--\(0.48\) to \(0.74\)--\(0.93\); per-family trader purity
climbs from \(0.09\)--\(0.14\) to \(0.71\)--\(0.77\). On average,
v4 catches \(83\%\) of manipulators at \(74\%\) purity, while the
baseline caught \(43\%\) at \(11\%\)---i.e.\ nine of every ten flags
were false positives. The single change was the node feature vector.

\begin{table}[t]
\caption{Within-generator OOD results, fifty runs.}
\label{tab:ood}
\centering
\begin{tabular}{lcc}
\toprule
Metric & v1 (2-feat) & v4 (8-feat) \\
\midrule
Edge AUC & 0.638 & \textbf{0.984} \\
Clique recall & 0.406 & \textbf{0.810} \\
Ring recall & 0.480 & \textbf{0.928} \\
Mixed recall & 0.408 & \textbf{0.745} \\
Clique purity & 0.113 & \textbf{0.774} \\
Ring purity & 0.136 & \textbf{0.729} \\
Mixed purity & 0.094 & \textbf{0.705} \\
\bottomrule
\end{tabular}
\end{table}

\subsection{Tier-2 Stacked Classifier}
Under leave-one-run-out cross-validation across the same fifty OOD
runs, the second-stage classifier consuming the baseline's score plus
the six engineered features achieves pooled trader-level AUC
\(0.968\)---a substantial lift from the baseline's \(0.793\)
trader-level AUC on the same data. At \(99\%\) precision the
classifier catches \(62\%\) of all manipulators in the cohort; at
\(95\%\) precision, \(74\%\); at the operating point of best F\(_1\),
\(76\%\) of manipulators at \(94\%\) purity. The lift over the
baseline (which caught \(11.7\%\) at the threshold typical of its own
score distribution) is roughly six-fold in recall at near-identical
purity.

We attribute the lift to three converging causes: feature stacking
(no single feature exceeds \(0.800\) standalone), non-linear
interaction modelling (a logistic regression on the same seven
features reaches AUC \(\approx 0.93\), so the boosted trees buy
roughly four more AUC points), and the differing error structure of
the baseline network versus the trader-marginal features. The
network's false positives tend to be high-activity benign traders; the
HHI feature's false positives tend to be traders with one or two
favourite counterparties; the overlap is mostly real manipulators.

\subsection{Feature Ablation Inside the Stacked Classifier}
\label{sec:ablate}
A leave-one-feature-out ablation through the stacked classifier
reveals an uneven distribution of load (Table~\ref{tab:abl}). Dropping
\(\phi_1\) (burst concentration)---the fifth-ranked feature by
standalone AUC---costs \(-0.051\) AUC, the largest single drop, and
collapses recall@99\%-purity from \(0.62\) to \(0.27\). Dropping
\(\phi_2\) (side entropy)---the worst feature standalone---costs only
\(-0.004\) AUC. The standalone ranking is therefore a poor proxy for
ensemble importance: features that are weak in isolation can be
load-bearing in combination.

\begin{table}[t]
\caption{Leave-one-feature-out ablation, tier-2 classifier.}
\label{tab:abl}
\centering
\begin{tabular}{lcc}
\toprule
Feature dropped & $\Delta$AUC & Rec.@p99 \\
\midrule
\textbf{full (7 features)} & --- & 0.619 \\
$\phi_1$ burst concentration & $-0.051$ & 0.274 \\
$\phi_6$ co-active count & $-0.007$ & 0.592 \\
$\phi_2$ side entropy & $-0.004$ & 0.634 \\
$\phi_4$ order qty CoV & $-0.003$ & 0.608 \\
$\phi_3$ counterparty HHI & $-0.002$ & 0.620 \\
$\phi_5$ top-partner share & $+0.000$ & 0.625 \\
$x_0$ baseline v1 score & $+0.000$ & 0.530 \\
\bottomrule
\end{tabular}
\end{table}

\subsection{Cross-Generator (ABIDES) Result}
\label{sec:abides}
Evaluating both detectors on the fifteen-run expanded ABIDES cohort
yields the comparison in Table~\ref{tab:abides}. The augmented network
retains edge AUC \(0.842\) on this genuinely unseen distribution; the
unaugmented baseline collapses to chance (\(0.518\)). The
\(+0.324\) AUC gap between the two models is consistent across all
runs (no fold sees the baseline above \(0.55\)). Per-family v4 mean
recall remains high on ABIDES: ring \(1.000\pm0.000\) across five
runs, mixed \(0.983\pm0.015\), and clique \(0.905\pm0.113\).

Purity, however, does \emph{not} transfer. Both detectors land in the
\(0.10\)--\(0.30\) purity range on ABIDES---far below the
\(0.70\)--\(0.77\) v4 achieves within-generator. The convergence of
both models on similar low purity suggests this is a property of the
\emph{cohort} (i.e., a locked threshold tuned in-sample on a smaller
cohort happens not to separate ABIDES benigns from manipulators
cleanly) rather than a v4-specific weakness. Cross-generator
deployment would therefore require per-generator threshold
recalibration, an honest finding that we surface explicitly rather
than paper over.

\begin{table}[t]
\caption{Cross-generator (ABIDES expanded, $n=15$) result.}
\label{tab:abides}
\centering
\begin{tabular}{lcc}
\toprule
Metric & v1 ABIDES & v4 ABIDES \\
\midrule
Edge AUC & 0.518 & \textbf{0.842} \\
Clique recall (mean$\pm$sd) & $0.867\pm0.16$ & $\mathbf{0.905\pm0.11}$ \\
Ring recall & $0.864\pm0.03$ & $\mathbf{1.000\pm0.00}$ \\
Mixed recall & $0.837\pm0.07$ & $\mathbf{0.983\pm0.02}$ \\
Clique purity (aggregate) & 0.096 & 0.094 \\
Ring purity & 0.276 & 0.293 \\
Mixed purity & 0.249 & 0.269 \\
\bottomrule
\end{tabular}
\end{table}

\subsection{Family-Disjoint Result}
\label{sec:family}
The most surprising result of the study comes from the
leave-one-family-out evaluation. For the stacked classifier of
\S\ref{sec:tier2} we trained three folds: one in which all clique runs
were excluded from training and the model was evaluated on the held-out
clique runs, and analogously for ring and mixed. The folds for the
end-to-end v4 network were generated identically, by adding a
\texttt{HOLDOUT\_FAMILY} filter to the trainer that excluded all runs of
that family from training. Each model was then evaluated on the full
fifty-run cohort.

Table~\ref{tab:family} contrasts the two architectures. The stacked
classifier retains AUC \(0.91\)--\(0.99\) on each held-out family,
with the ring fold actually exceeding the baseline leave-one-run-out
reference (\(0.994\) vs.\ \(0.968\)). The end-to-end network, by
contrast, collapses to chance or anti-chance on every fold.

\begin{table}[t]
\caption{Family-disjoint AUC. Tier-2 is leave-one-family-out CV; v4 is end-to-end retraining with the family filtered from training data.}
\label{tab:family}
\centering
\begin{tabular}{lcc}
\toprule
Held-out family & Tier-2 GBM AUC & v4 GraphSAGE AUC \\
\midrule
Ring & \textbf{0.994} & 0.480 \\
Clique & 0.975 & 0.533 \\
Mixed & 0.906 & 0.374 \\
\bottomrule
\end{tabular}
\end{table}

A methodological caveat applies to the v4 result. Excluding a family
from training reduces each per-day cohort from approximately fifty runs
to thirty-three, and the validation-set sampling correspondingly
shrinks from five runs to three. With only three validation graphs,
the early-stopping signal (mean per-family trader recall on validation)
saturates after one or two epochs, terminating training at eight
epochs of overfitting on a tiny in-family sample. A longer training
schedule with looser early stopping might recover some performance and
remains worth investigating. Nevertheless, the contrast with the
stacked classifier on the same task is unambiguous: under the standard
training protocol that produced the headline v4 OOD result, family
disjointness is fatal to the end-to-end network.


\subsection{Per-Run Variance on the ABIDES Cohort}
\label{sec:abides-var}
The per-run breakdown of the cross-generator result deserves attention
because three of the fifteen ABIDES runs are responsible for most of
the spread in clique recall. Across the five ring runs, v4 recall is
exactly \(1.000\) with zero variance; across the five mixed runs,
recall is \(0.98\pm0.02\); across the five clique runs, recall is
\(0.91\pm0.11\) with one outlier at \(0.69\). The outlier corresponds
to a clique scenario whose front counterparties happened to receive
unusually balanced order flow, partially masking the burst-concentration
signature. The remaining four clique runs all exceed \(0.85\). This
pattern, where the same metric is bimodal across runs, recurs in the
within-generator OOD cohort as well, and is a fair argument for
reporting cohort-mean numbers alongside per-run distributions.

\subsection{Case Study: A Successful Ring Detection}
\label{sec:case}
We pick a representative ring run from the within-generator OOD cohort
and trace why v4 succeeds. The run contains \(\sim\!200\) traders, six
of whom are members of a ring rotating positions among themselves over
a single \(\sim\!45\)-minute window. The six manipulator traders share
three salient properties:
(i) extreme burst concentration \(\phi_1 > 0.85\) (all six)---versus
a benign median of \(0.16\);
(ii) low counterparty HHI \(\phi_3 < 0.2\), reflecting the ring's
intentionally diversified counterparty list, versus a benign median
near \(0.5\) (note the manipulator direction is \emph{low}, since
genuinely focused traders tend to have a small partner set, the
opposite of what the textbook intuition would predict);
(iii) high co-active count \(\phi_6\) because the six rotate trades
among themselves and against each other simultaneously.
The augmented network's first SAGEConv layer absorbs these node-level
signals, and the message-passing step pools each ring member's
neighbourhood (which is the rest of the ring), reinforcing the
classification. The same six traders are flagged by the second-stage
classifier independently on the trader-marginal features alone.

\subsection{What Happens When v4 Fails}
\label{sec:fail}
In the family-disjoint regime, v4's early-stopping signal hits a
plateau-like maximum within two epochs. This is unsurprising once the
val set is reduced to three graphs: with \(\sim\!1{,}200\) trader nodes
total and a handful of positives, any non-trivial scoring function will
achieve perfect trader recall on the val set at some threshold. The
patience-\(6\) counter then expires without progress, and training
terminates at epoch \(8\). The remaining \(75\%\) of the epoch budget
is unused. We believe this is the primary reason the v4 family-disjoint
result collapsed, and that a longer schedule or larger val pool may
substantially change the picture.

\section{Discussion}
\label{sec:disc}

The pattern of results across these three out-of-distribution regimes
suggests a clear architectural distinction. The engineered
manipulation-signature features carry signal that is
\emph{manipulation-essential rather than family-specific}: tier-2 on
those features generalises across families well above the
all-family leave-one-run-out reference. Yet the end-to-end network,
which consumes the same six features as node attributes, does not
generalise across families. The most natural explanation is that the
GNN's message-passing layers learn to recognise family-specific
structural neighbourhoods (a ring's closed cycle is structurally
different from a clique's shared-front arrangement), and when one
family is held out the network's structural prior simply has no
correspondence to the held-out pattern.

The bolt-on classifier in contrast sees each trader independently, with
no notion of structural neighbourhood at all---it cannot use the
neighbourhood signal, but it cannot be misled by it either. The
trader-marginal features are family-agnostic by construction
(burst concentration, side entropy, etc.\ describe what manipulation
looks like, not what clique manipulation looks like), and the GBM
learns combinations that transfer.

This is the most actionable finding of the study. For deployments
where the operational distribution may differ structurally from the
training distribution---a different exchange, a different ticker
universe, a different manipulation taxonomy than was anticipated at
training time---a simpler classifier on top of taxonomy-grounded
features may be substantially more robust than a more expressive
end-to-end network. The latter wins comfortably within-distribution,
but its advantage is concentrated in regimes where the training and
test distributions share structural priors.

The cross-generator (ABIDES) result occupies an interesting middle
ground: the engineered features generalise (v4 holds AUC \(0.842\)
versus the baseline's chance) but the operating point does not (purity
collapses uniformly). This is consistent with the interpretation that
\emph{ranking} is preserved across the cross-generator shift but the
absolute manipulator score distribution is shifted enough that a
threshold tuned on one cohort flags too many benigns on the other. In
production, this argues for per-cohort threshold calibration as part of
deployment hygiene rather than for retraining at every shift.

\section{Limitations and Future Work}
\label{sec:lim}

Several caveats accompany the present results. First, both the primary
generator and ABIDES are synthetic; no real-data labelled cohort is
involved, and although the synthesiser is calibrated to NSE bhavcopy
moments, real markets contain heavy-tailed events and microstructure
shifts that synthetic data captures only approximately. Second, the
ABIDES cohort, although now at \(n=15\), remains modest; expanding to
fifty or more runs would tighten the confidence intervals on the
cross-generator finding. Third, the v4 family-disjoint result was
obtained under the standard training schedule of the main OOD run, and
a separate study with a longer schedule and larger validation set is
indicated before the architectural conclusion is locked in.

Future work falls into three threads. The most direct is a real-data
pilot, in which the detectors are run on labelled real-market data
(where labels can be obtained from regulatory enforcement actions); the
absence of labelled real manipulation data has been a recurring
limitation. Second, the family-disjoint v4 retraining should be redone
with a larger validation pool (e.g.\ \(20\%\) per day) and longer
patience (e.g.\ \(20\) epochs), to determine whether v4's collapse is
intrinsic or an early-stopping artefact. Third, end-to-end models may
benefit from family-agnostic structural priors---e.g.\ randomly
permuting node features between families during training---to force
the message-passing layers to be less family-specific.


\section{Threats to Validity}
\label{sec:threats}
We highlight four threats to the generalisability of the conclusions
above. First, the locked decision threshold for both detectors is
tuned in-sample on the evaluation cohort, which inflates absolute
purity numbers; the cross-MODEL comparison is fair (both detectors
suffer the same way), but the absolute operating point is mildly
optimistic. Second, the cross-generator cohort, while drawn from a
genuinely different simulator, is still synthetic. Real-market
manipulators may exhibit microstructure patterns absent from both
synthesisers. Third, the family-disjoint v4 retraining used the
training schedule tuned for the original five-day cohort, and the
schedule's interaction with a reduced cohort and reduced validation
pool is responsible for at least some of the observed collapse. A
follow-up retraining with a relaxed early-stopping criterion is
indicated. Fourth, the engineered features were designed by hand based
on a taxonomy of manipulation patterns, which is itself a form of
inductive bias; features tailored to a different taxonomy would yield
a different decision surface.

\section{Related Threads Not Pursued}
\label{sec:notpursued}
For transparency we list two directions we did not pursue here but
which would naturally extend the work. The first is a benign-only
cohort, in which every run contains no manipulators; running either
detector against such a cohort would yield a clean false-positive rate
estimate at any chosen threshold and is the proper way to report the
``cost'' side of a flagging policy. The second is a study of detector
behaviour against an adaptive adversary, in which the manipulator
agents themselves are trained or hand-tuned to evade the detector. We
note both directions but defer them to subsequent work.

\section{Conclusion}
\label{sec:concl}

We have presented a multi-cohort study of a two-layer GraphSAGE
detector augmented with six manipulation-signature features and
compared it against a stacked second-stage classifier consuming the
same features. The augmented network achieves edge AUC \(0.984\) and
recall \(0.83\) at purity \(0.74\) on a fifty-run within-generator OOD
cohort, lifting from a two-feature baseline of \(0.638\)/\(0.43\)/\(0.11\).
The stacked classifier achieves pooled AUC \(0.968\) with \(62\%\)
recall at \(99\%\) purity. Both detectors are then evaluated on a
fifteen-run cross-generator cohort drawn from a different agent-based
simulator: the augmented network retains AUC \(0.842\) versus the
baseline's chance, demonstrating that the engineered features carry
signal across generators. A leave-one-family-out evaluation reveals an
unexpected contrast, with the stacked classifier retaining AUC
\(0.91\)--\(0.99\) on each held-out family while the end-to-end
network collapses to chance. The simpler, trader-marginal classifier
therefore generalises across manipulation styles better than the more
expressive end-to-end network. We argue that for surveillance settings
where the operational distribution may differ structurally from the
training distribution, the bolt-on architecture deserves serious
consideration as a deployment choice on robustness grounds, even when
the end-to-end model wins within-distribution comparisons.

\begin{thebibliography}{99}

% Post-audit bibliography (v2, 2026-08-06). Ten entries; all high-confidence
% real and independently DOI-verifiable. Ten problematic entries removed
% during Workstream A of the pre-submission audit remediation:
%   khanna2021review  — journal "Journal of Financial Surveillance" not
%                        verifiable in any Springer/ACM/IEEE index
%   sun2022fraudgnn   — key/content correspondence unverifiable
%   cong2020amlsim    — attribution wrong (AMLSim is Suzumura & Kanezashi)
%   nguyen2023iclr    — cited ICLR workshop "ML for Finance" does not exist
%   leskovec2014snap  — key year 2014 vs entry year 2016 mismatch
%   chen2016xgboost, wang2019fdgars, liu2021pickfeat, xu2019gin,
%   dou2020carefraud — uncited in v1; will be re-added in Workstream E1
%                       when concrete external-baseline comparisons land
% See docs/AUDIT.md § B10-B11 and docs/REMEDIATION_PLAN.md workstream A.

\bibitem{cao2014pso}
Y. Cao, Y. Li, S. Coleman, A. Belatreche, and T. McGinnity,
``Particle-swarm-optimised support-vector machine for detecting
market manipulation,'' in \emph{Proc. IEEE Symp. on Computational
Intelligence for Financial Engineering}, 2014.

\bibitem{hamilton2017graphsage}
W. L. Hamilton, R. Ying, and J. Leskovec, ``Inductive representation
learning on large graphs,'' in \emph{Advances in Neural Information
Processing Systems (NeurIPS)}, 2017.

\bibitem{velickovic2018gat}
P. Veli\v{c}kovi\'{c}, G. Cucurull, A. Casanova, A. Romero, P. Li\`{o},
and Y. Bengio, ``Graph attention networks,'' in \emph{Int. Conf. on
Learning Representations (ICLR)}, 2018.

\bibitem{kipf2017gcn}
T. N. Kipf and M. Welling, ``Semi-supervised classification with graph
convolutional networks,'' in \emph{Int. Conf. on Learning
Representations (ICLR)}, 2017.

\bibitem{lin2017focal}
T.-Y. Lin, P. Goyal, R. Girshick, K. He, and P. Doll\'{a}r,
``Focal loss for dense object detection,'' in \emph{Proc. IEEE Int.
Conf. on Computer Vision (ICCV)}, 2017.

\bibitem{friedman2001gradient}
J. H. Friedman, ``Greedy function approximation: a gradient boosting
machine,'' \emph{Annals of Statistics}, vol. 29, no. 5, pp. 1189--1232,
2001.

\bibitem{recht2019imagenet}
B. Recht, R. Roelofs, L. Schmidt, and V. Shankar, ``Do ImageNet
classifiers generalize to ImageNet?'' in \emph{Proc. Int. Conf. on
Machine Learning (ICML)}, 2019.

\bibitem{koh2021wilds}
P. W. Koh, S. Sagawa, H. Marklund, S. M. Xie, M. Zhang, A. Balsubramani,
W. Hu, M. Yasunaga, R. L. Phillips, I. Gao, T. Lee, E. David, I. Stavness,
W. Guo, B. A. Earnshaw, I. S. Haque, S. Beery, J. Leskovec, A. Kundaje,
E. Pierson, S. Wang, D. Koller, S. Levine, and P. Liang, ``WILDS: A
benchmark of in-the-wild distribution shifts,'' in \emph{Proc. Int.
Conf. on Machine Learning (ICML)}, 2021.

\bibitem{byrd2019abides}
D. Byrd, M. Hybinette, and T. H. Balch, ``ABIDES: Towards
high-fidelity multi-agent market simulation,'' in \emph{Proc. ACM
SIGSIM Conf. on Principles of Advanced Discrete Simulation}, 2019.

\bibitem{cont2001stylized}
R. Cont, ``Empirical properties of asset returns: stylized facts and
statistical implications,'' \emph{Quantitative Finance}, vol. 1,
no. 2, pp. 223--236, 2001.

\end{thebibliography}

\end{document}
